%------------------------------------------------------------------------------------
%  SECTION 3 : GESTION DES FIGURES ET ILLUSTRATIONS
%------------------------------------------------------------------------------------
\section{Gestion des Figures et Illustrations}
\label{sec:illustrations}

Le modèle sépare les figures techniques conventionnelles des illustrations visuelles ou artistiques grâce à une liste dédiée.

\subsection{Figure standard avec sous-figures}
\label{subsec:figures_standards}

La \cref{fig:exemple_figures} présente une organisation de sous-figures alignées avec des légendes indépendantes.

\begin{figure}[htbp]
  \centering
  \begin{subfigure}[b]{0.45\textwidth}
    \centering
    \rule{0.8\textwidth}{3cm}
    \caption{Schéma technique A}
    \label{subfig:schema_a}
  \end{subfigure}
  \hfill
  \begin{subfigure}[b]{0.45\textwidth}
    \centering
    \rule{0.8\textwidth}{3cm}
    \caption{Schéma technique B}
    \label{subfig:schema_b}
  \end{subfigure}
  \caption{Exemple de figure double avec sous-légendes}
  \label{fig:exemple_figures}
\end{figure}

\subsection{Environnement Illustration personnalisé avec cadre décoratif}
\label{subsec:illustrations_dediees}

L'environnement \texttt{illustration} permet d'indexer les supports iconographiques dans une table distincte.
Le cadre personnalisé \texttt{cadreFigure} permet d'englober l'image d'un contour soigné aux coins arrondis :

\begin{illustration}[htbp]
  \centering
  \begin{cadreFigure}[width=0.7\textwidth]
    \rule{7cm}{3.5cm}
  \end{cadreFigure}
  \caption{Représentation anatomique ou graphique d'expérimentation}
  \label{ill:anatomie}
\end{illustration}

L'\cref{ill:anatomie} apparaît automatiquement dans la \emph{Liste des illustrations}.
